% !TeX root = ../qm_chap3.tex

\begin{frame}{教材3.7　从重复测量的频率得到期望值}
  \iqblayouttwo{
    \iqbsectiontitle{离散测量的平均值}
    \begin{iqbitemize}
      \item 对$N$个采用同一方式制备的独立系统分别测一次。若结果$b$出现$n_b$次，样本均值是$\sum_b b\,n_b/N$。
      \item 样本足够多时，相对频率$n_b/N$趋近概率$P_b$，均值趋近$\sum_b bP_b$。
      \item 例如等概率测得$0$和$2$时，期望值为$1$，它可以不在单次测量的候选结果中。
    \end{iqbitemize}
  }{
    \iqbsectiontitle{连续位置与一般算符}
    \begin{iqbitemize}
      \item 求三维系统中$x$的分布时，对$y,z$积分得到边缘概率密度。再乘以$x$并积分，得到$\langle x\rangle$。
      \item 用归一化态计算一般物理量的期望值：
      \[
        \langle B\rangle=\int\Psi^*(\hat B\Psi)\,\mathrm{d}\tau .
      \]
      \item 期望值对算符加法和数乘是线性的。即使$A=B$，也可能有$\langle A^2\rangle\ne\langle A\rangle^2$。
    \end{iqbitemize}
  }
\end{frame}

\begin{frame}{教材3.7　三维箱中基态的平均位置在箱中心}
  \iqbfontsizemode{small}
  \iqbsectiontitle{计算位置期望值}

  \textbf{问题}：三维箱中粒子处于基态$\psi_{111}=X(x)Y(y)Z(z)$，取$x$方向的箱宽为$l=a$。

  \begin{iqbitemize}
    \item $Y$和$Z$各自归一化，计算$x$的期望值时，对$y,z$的积分各给出1，只剩$x$方向的基态因子：
    \[
      X(x) = \sqrt{\dfrac{2}{l}}\sin\left(\dfrac{\pi x}{l}\right), \quad 0 \leq x \leq l
    \]

    \item \textbf{坐标平均值}：
    \begin{align*}
      \langle x \rangle &= \int_0^l x|X(x)|^2\,\mathrm{d}x \\
      &= \dfrac{2}{l}\int_0^l x\sin^2\left(\dfrac{\pi x}{l}\right)\,\mathrm{d}x \\
      &= \dfrac{2}{l} \cdot \dfrac{l^2}{4} = \dfrac{l}{2}
    \end{align*}
    $x$的概率密度关于$l/2$对称，且在该点达到峰值。
  \end{iqbitemize}

  \iqborangebox{三个方向}{对$y,z$同理，三维基态的平均位置为$(a/2,b/2,c/2)$。}
\end{frame}

\begin{frame}{箱宽决定基态的位置标准差}
  \iqbfontsizemode{small}
  \iqbsectiontitle{计算位置方差}

  \begin{iqbitemize}
    \item \textbf{坐标平方平均值}：
    \begin{align*}
      \langle x^2 \rangle &= \dfrac{2}{l}\int_0^l x^2\sin^2\left(\dfrac{\pi x}{l}\right)dx \\
      &= l^2\left(\dfrac{1}{3} - \dfrac{1}{2\pi^2}\right)
    \end{align*}

    \item \textbf{位置不确定度}：
    \[
      (\Delta x)^2 = \langle x^2 \rangle - \langle x \rangle^2 = l^2\left(\dfrac{1}{12} - \dfrac{1}{2\pi^2}\right)
    \]

    \item \textbf{标准差}：
    \[
      \Delta x = \dfrac{l}{2}\sqrt{\dfrac{1}{3} - \dfrac{1}{2\pi^2}}
    \]
  \end{iqbitemize}

  \iqbgreenbox{结论}{位置不确定度与箱宽$l$成正比}
\end{frame}

\begin{frame}{教材3.7　平均动量为零，动量方差仍为正}
  \iqbfontsizemode{small}
  \iqbsectiontitle{计算动量期望值}

  \begin{iqbitemize}
    \item $Y$和$Z$的积分仍各等于1。$X$为实函数，在两端取零，因此
    \[
      \langle p_x\rangle=-i\hbar\int_0^l XX'\,\mathrm{d}x
        =-\dfrac{i\hbar}{2}[X^2]_0^l=0 .
    \]

    \item \textbf{动量平方平均值}：
    \[
      \langle p_x^2 \rangle = \dfrac{2\hbar^2\pi^2}{l^3}\int_0^l \sin^2\left(\dfrac{\pi x}{l}\right)dx = \dfrac{h^2}{4l^2}
    \]

    \item \textbf{动量不确定度}：
    \[
      \Delta p_x = \sqrt{\langle p_x^2 \rangle - \langle p_x \rangle^2} = \dfrac{h}{2l}
    \]

    \item \textbf{验证不确定性原理}：
    \[
      \Delta x \cdot \Delta p_x = \dfrac{h}{2}\sqrt{\dfrac{1}{12} - \dfrac{1}{2\pi^2}} > \dfrac{\hbar}{2}
    \]
  \end{iqbitemize}

  \iqbgreenbox{核心}{虽然$\langle p_x \rangle = 0$，但$\langle p_x^2 \rangle > 0$体现零点能}
\end{frame}

\begin{frame}{微扰理论从可精确求解的模型出发}
  \iqblayouttwo{
  
    \iqbsectiontitle{微扰理论的基本思想}
    \begin{iqbitemize}
    \item \textbf{核心思想}（perturbation theory）：
    \subitem{大多数量子系统无法精确求解}
    \subitem{将复杂哈密顿分为可解部分 + 微扰部分}
    \subitem{通过级数展开获得近似解}
    \end{iqbitemize}

    
  }{
  
    \iqbsectiontitle{哈密顿算符的分解}
    \begin{iqbitemize}
    \item \textbf{分解方案}：
    \end{iqbitemize}
    \[
      \hat{H} = \hat{H}_0 + \lambda\hat{H}'
    \]
    \begin{iqbitemize}
    \item \textbf{各项含义}：
    \subitem{$\hat{H}_0$：未微扰哈密顿（可精确求解）}
    \subitem{$\hat{H}'$：微扰项（扰动）}
    \subitem{$\lambda$：小参数，标记微扰强度}
    \item \textbf{已知条件}：
    \subitem{$\hat{H}_0\phi_n^{(0)} = E_n^{(0)}\phi_n^{(0)}$ 已知}
    \subitem{目标是求$\hat{H}$的近似本征值$E_n$和本征函数$\psi_n$}
    \end{iqbitemize}

    \iqborangebox{策略}{将难解问题转化为可解问题 + 小修正的级数}
  
  }
\end{frame}

\begin{frame}{一阶能量修正等于微扰在零阶态的期望值}

  
    \iqbsectiontitle{非简并一阶微扰公式}
    \begin{iqbitemize}
    \item \textbf{能量的一阶修正}：
    \end{iqbitemize}
    \[
      E_n \approx E_n^{(0)} + \langle\phi_n^{(0)}|\hat{H}'|\phi_n^{(0)}\rangle
    \]
    \begin{iqbitemize}
    \item \textbf{物理意义}：
    \subitem{能量修正 = 微扰在对角态上的期望值}
    \subitem{仅涉及单个态$n$，计算相对简单}
    \end{iqbitemize}

    \iqbsep

    \end{frame}

\begin{frame}{微扰混入其他本征态并修正波函数}
  \iqblayouttwo{
  
    \iqbsectiontitle{波函数的一阶修正}
    \begin{iqbitemize}
    \item \textbf{修正公式}：
    \end{iqbitemize}
    \[
      \psi_n \approx \phi_n^{(0)} + \sum_{k \neq n}\dfrac{\langle\phi_k^{(0)}|\hat{H}'|\phi_n^{(0)}\rangle}{E_n^{(0)} - E_k^{(0)}}\phi_k^{(0)}
    \]
    \begin{iqbitemize}
    \item \textbf{物理意义}：
    \subitem{微扰使态$n$与其他所有态$k \neq n$混合}
    \subitem{混合强度由矩阵元$\langle k|\hat{H}'|n\rangle$和能级间距$E_n^{(0)} - E_k^{(0)}$决定}
    \subitem{能级越近，耦合越强}
    \end{iqbitemize}

    \iqbsep

    
  }{
  
    \iqbsectiontitle{微扰理论的适用性}
    \begin{iqbitemize}
    \item \textbf{适用条件}：
    \subitem{微扰项$\lambda\hat{H}'$远小于未微扰项}
    \subitem{级数收敛，通常只需一阶或二阶}
    \item \textbf{不适用情况}：
    \subitem{简并能级（需用简并微扰理论）}
    \subitem{强耦合系统}
    \end{iqbitemize}

    \iqbgreenbox{核心}{能量修正是对角矩阵元，波函数修正来自所有其他态的混合}
  
  }
\end{frame}

\begin{frame}{在谐振子势能中加入四次项}
  \iqblayouttwo{
  
    \iqbsectiontitle{实例：一维非谐振子}
    \begin{iqbitemize}
    \item \textbf{物理系统}：
    \subitem{实际分子振动不是完全谐性的}
    \subitem{大位移时势能偏离抛物线}
    \subitem{需加入高阶项修正}
    \end{iqbitemize}

    
  }{
  
    \iqbsectiontitle{哈密顿算符}
    \begin{iqbitemize}
    \item \textbf{完整哈密顿}：
    \end{iqbitemize}
    \[
      \hat{H} = -\dfrac{\hbar^2}{2m}\dfrac{\mathrm{d}^2}{\mathrm{d}x^2} + \dfrac{1}{2}kx^2 + \lambda x^4
    \]
    \begin{iqbitemize}
    \item \textbf{分解方案}：
    \subitem{$\hat{H}_0 = -\dfrac{\hbar^2}{2m}\dfrac{d^2}{dx^2} + \dfrac{1}{2}kx^2$（谐振子，已知解）}
    \subitem{$\hat{H}' = x^4$（四次项微扰）}
    \subitem{$\lambda$：小参数，控制非谐性强度}
    \end{iqbitemize}

    \iqbsmallsep

    
  }
\end{frame}

\begin{frame}{四次项改变非谐振子的能级}
  \iqblayouttwo{
  
    \iqbsectiontitle{能量修正的计算}
    \begin{iqbitemize}
    \item \textbf{一阶能量修正}：
    \subitem{计算矩阵元：$\langle n|x^4|n\rangle$}
    \subitem{利用升降算符：$x \propto (\hat{a}_+ + \hat{a}_-)$}
    \subitem{结果：$E_n^{(1)} = \dfrac{3\hbar^2}{4m^2\omega^2}(2n^2 + 2n + 1)\lambda$}
    \item \textbf{物理意义}：
    \subitem{能级不再是等间距的}
    \subitem{高激发态非谐性更显著}
    \end{iqbitemize}

    \iqborangebox{实例}{分子振动的非谐性解释了红外光谱的精细结构}
  
  }{
  
    \iqbsectiontitle{微扰理论的实际应用}

    
  }
\end{frame}

\begin{frame}{外场微扰导致原子谱线分裂或移动}

  
    \iqbsectiontitle{原子光谱效应}
    \begin{iqbitemize}
    \item \textbf{斯塔克效应}（Stark effect）：
    \subitem{外电场$\mathbf{E}$对原子能级的影响}
    \subitem{微扰项：$\hat{H}' = -\boldsymbol{\mu} \cdot \mathbf{E}$}
    \subitem{能级分裂，谱线分裂}
    \subitem{一阶斯塔克效应（氢原子）与二阶斯塔克效应}

    \item \textbf{塞曼效应}（Zeeman effect）：
    \subitem{外磁场$\mathbf{B}$对原子能级的影响}
    \subitem{微扰项：$\hat{H}' = -\boldsymbol{\mu}_L \cdot \mathbf{B}$}
    \subitem{能级按磁量子数$m_\ell$分裂}
    \subitem{正常塞曼效应与反常塞曼效应}
    \end{iqbitemize}

    \iqbsep

    \end{frame}

\begin{frame}{微扰理论修正量子化学的近似结果}
  \iqblayouttwo{
  
    \iqbsectiontitle{化学中的微扰理论}
    \begin{iqbitemize}
    \item \textbf{分子间相互作用}：
    \subitem{范德华力：偶极-偶极相互作用}
    \subitem{氢键：静电能 + 诱导能 + 色散能}
    \subitem{用微扰理论计算结合能}

    \item \textbf{化学反应理论}：
    \subitem{前线轨道理论：HOMO-LUMO相互作用}
    \subitem{微扰处理轨道混合}
    \subitem{预测反应活性位点和方向}

    \item \textbf{晶体场理论}：
    \subitem{配体场对金属$d$轨道的微扰}
    \subitem{解释配合物的颜色和磁性}
    \end{iqbitemize}
  }{
  
    \iqbsectiontitle{量子化学计算方法}
    \begin{iqbitemize}
    \item \textbf{多体微扰理论}（MBPT）：
    \subitem{MP2、MP3、MP4等方法}
    \subitem{考虑电子相关}
    \item \textbf{组态相互作用}（CI）：
    \subitem{激发态的混合}
    \subitem{高精度电子结构计算}
    \end{iqbitemize}

    \iqbgreenbox{重要性}{微扰理论是量子化学计算的基础，从简单模型到高精度计算}
  
  }
\end{frame}
